\section{Fixtures}

The positive fixture is \texttt{mail-batch-added}. It represents a canvassed
state where precinct P-001 has an election-day batch and a late mail batch.
The status event explains that the late mail batch was accepted before canvass.

\begin{center}
\small
\begin{tabularx}{\textwidth}{lrrr}
\toprule
Batch & Accepted & Counted & Rejected \\
\midrule
\texttt{batch:P-001:election-day} & 70 & 70 & 0 \\
\texttt{batch:P-001:late-mail} & 10 & 10 & 0 \\
\texttt{batch:P-002:election-day} & 60 & 60 & 0 \\
\bottomrule
\end{tabularx}
\end{center}

The fixture uses these anchors:

\begin{center}
\small
\begin{tabularx}{\textwidth}{lX}
\toprule
Path & Role \\
\midrule
\path{normalized/batches.ndjson} & Declares batch manifest rows. \\
\path{normalized/summaries.ndjson} & Declares batch-level and jurisdiction summaries. \\
\path{status/events.ndjson} & Explains the late-mail transition into the canvassed state. \\
\path{transcripts/verify-transcript.json} & Records verifier pass/fail checks. \\
\bottomrule
\end{tabularx}
\end{center}

The negative fixture is \texttt{missing-batch}. Its summaries still reference
the P-001 late-mail batch, but \path{normalized/batches.ndjson} omits that
manifest row. Source hashes and package structure can still be present; the
accounting evidence is incomplete.
